A three-dimensional fluid can stretch its own rotation.
Take a region where nearby fluid rotates around a common axis.
If the velocity field lengthens that region along the axis, incompressibility forces its cross-section to narrow.
When the rotation remains aligned with the extensional direction, the strain amplifies the vorticity as the region narrows.
If that alignment survives, the thinner and stronger region can be stretched again.
That is the vortex we have to deal with: the motion can repeat.

The equation keeps the stretching and the viscous response inside that one motion.
Let
\[
  \omega:=\nabla\times u,
  \qquad
  S:=\frac12\bigl(\nabla u+(\nabla u)^{\mathsf T}\bigr).
\]
Taking the curl of \NavierStokes{} gives
\begin{equation}\label{eq:VorticityEquation}
  \partial_t\omega+(u\cdot\nabla)\omega
  =S\omega+\nu\Delta\omega.
\end{equation}
The left side follows the vorticity with the moving fluid.
The term \(S\omega\) is the aligned stretch we just described.
The term \(\nu\Delta\omega\) is the viscous response to the spatial profile being sharpened as the vortex narrows.
The contraction changes both terms at once.

At a smaller width, the decisive comparison is the relative scale of \(S\omega\) and \(\nu\Delta\omega\).
The \NavierStokes{} rescaling computes it.
